% !TEX root = ../../../main.tex
% Appendix entry: detection.sax_dsax -- "Symbolic encoding (dSAX)" (Appendix A).
\entryfield{Algorithm}
\code{_run} in \code{Adapters/detection_sax_dsax.py} resolves the segment length with the
shared \code{segment_plan} (see \blockentry{detection.sax_csax}) and calls \code{dsax} in
\code{Working/Detection/sax/dsax_python/dsax.py}, with the whole span as training set. An
\emph{absolute\_threshold} or \emph{min\_same\_halfwidth} of 0 is passed on as ``unset''.
With $m$ segments of $s$ samples and alphabet size $k$:
\begin{enumerate}
\item \emph{Checks.} Unknown estimator or mode, $k<2$, $m<2$ or $s<2$ raise
  \code{ValueError}; so does \emph{absolute} mode with $k\neq 3$ or no threshold, and
  \emph{quantile} mode with \emph{same\_fraction} outside $(0,1)$.
\item \emph{Trim} to $n' = ms$ samples. Mark every non-finite sample as missing; a segment
  that touches one is a \emph{missing segment}.
\item \emph{Normalise} once (unlike cSAX and pSAX): with \emph{normalize} on, subtract the
  mean $\mu$ of the recorded samples and divide by their standard deviation $\sigma$, unless
  $\sigma < 0.001$ in the stored unit, in which case only subtract.
\item \emph{Measure each segment's rise} $\delta_i$ (\code{compute_deltas} in
  \code{trend_estimators.py}), always as rise across the segment, so estimators are
  interchangeable:
  \begin{itemize}
  \item \code{endpoints}: last sample minus first;
  \item \code{robust_endpoints}: median of the last $k_e$ samples minus median of the first
    $k_e$, $k_e = \min(\text{endpoint\_k}, \lfloor s/2\rfloor)$;
  \item \code{ols_slope} (default): least-squares slope against sample index, times $s-1$;
  \item \code{theil_sen}: \code{scipy.stats.theilslopes} slope, times $s-1$, on at most 200
    evenly spaced samples per segment.
  \end{itemize}
  Missing segments are measured on a zero-filled copy and their $\delta_i$ then set to NaN.
\item \emph{Cutlines} in the working delta domain, from the recorded segments' deltas:
  \begin{itemize}
  \item all deltas zero (within $10^{-9}$ of the signal's scale), in \emph{learned} or
    \emph{quantile} mode: a band wide enough that every segment is the middle letter;
  \item all deltas equal but not zero (a clean ramp): equal-width cutlines across
    $[-|v|, |v|]$, so the ramp reads all-up or all-down;
  \item \emph{learned}: pSAX's quantiser on the deltas, unchanged (Epanechnikov density,
    k-means++ start with the fixed seed, Lloyd--Max; see \blockentry{detection.sax_psax});
  \item \emph{absolute}: $\pm\,\text{absolute\_threshold}/\sigma$ (the threshold is in raw
    units of rise per segment);
  \item \emph{quantile}: for $k=3$ the $\tfrac{1-f}{2}$ and $\tfrac{1+f}{2}$ quantiles with
    $f$ = \emph{same\_fraction}; for $k>3$ equiprobable quantiles $j/k$.
  \end{itemize}
  With \emph{force\_symmetric} on (learned and quantile modes), fold the cutlines about zero:
  $c'_j = (c_j - c_{k-2-j})/2$. Then widen the middle band to at least
  $\pm$\emph{min\_same\_halfwidth} if one is given (odd $k$ only), and nudge any cutlines
  that collide apart by $10^{-9}$ of their scale.
\item \emph{Quantise}: $\text{symbol}_i = \#\{j : c_j \leq \delta_i\}$
  (\code{np.searchsorted}, right side); missing segments get $-1$.
\end{enumerate}
The live read-out (\code{derive_sax_rows}) and the declared but unused span-aware defaults
are the shared ones described under \blockentry{detection.sax_csax}.

\emph{The 14 settings, briefly.} Five set the segment length (\emph{segment\_mode} and the
four controls it chooses between). \emph{alphabet\_size} (2--9) and \emph{trend\_estimator}
set what is measured and into how many letters; \emph{endpoint\_k} matters only for
\code{robust_endpoints}. \emph{threshold\_mode} chooses the cutline rule, and
\emph{absolute\_threshold} and \emph{same\_fraction} feed only the \emph{absolute} and
\emph{quantile} modes. \emph{min\_same\_halfwidth} and \emph{force\_symmetric} adjust the
cutlines afterwards, and \emph{normalize} chooses the working domain.

\entryfield{Assumptions}
\begin{itemize}
\item \emph{Units.} Not unit-aware. \emph{absolute\_threshold} is in the stored unit per
  segment: the same number means a thousand times more on a recording stored in volts than
  on one stored in millivolts. \emph{min\_same\_halfwidth} is in the \emph{working} domain,
  that is, in standard deviations of the span when \emph{normalize} is on. The fixed
  $0.001$ normalisation threshold is in the stored unit.
\item \emph{Missing values} (since 2026-10-14). A segment touching a NaN or infinity is
  symbol $-1$, spelt \emph{-} by the bridge and by \blockref{detection.symbol_search}. The
  normalisation (which used to be the third statistic a single NaN poisoned), the cutlines,
  the observed ``same'' share and the mean delta are all taken over the recorded samples or
  segments. A span with no recorded segment comes out all $-1$. Gap-free input gives the
  same result as before the change, bit for bit. Unlike the five-stage encoding, the
  read-out gives no \emph{Missing samples} warning before a run.
\item \emph{Baseline.} A trend encoder reads drift as morphology; remove the baseline
  first. The shipped template pairs it with baseline removal
  (\code{webui/server/templates.py}).
\item \emph{Sampling rate.} The rise is per segment, not per second, so at matched seconds
  per symbol the encoding does not depend on the sampling rate (test 11 below).
\end{itemize}

\entryfield{Outputs}
An \code{Encoding} of kind \code{symbolic}: $m$ symbols in $0,\dots,k-1$ (0 the steepest
fall) or $-1$. \code{meta["details"]} holds the shared keys (segment length and count,
trimming, \code{paa}, \code{paa_raw}, cutlines, representatives, normalisation constants)
and dSAX's own: \code{cutline_domain} \code{= "delta"}, \code{deltas}, \code{deltas_raw},
\code{delta_scale} ($\sigma$ only: a difference takes no offset), \code{delta_mean},
\code{seg_slope}, \code{segment_starts}, \code{segment_ends}, \code{zero_symbol},
\code{same_fraction_observed}, \code{cutlines_degenerate}, \code{theil_sen_subsampled},
\code{min_same_halfwidth_applied} and \code{n_missing_segments}. The bridge sends the
symbols spelt \emph{a, b, c}\ldots{} (not the D/S/U of \code{dsax_letters}) with the
working-domain cutlines and their domain.

\entryfield{Cost}
The three default estimators are vectorised, $O(n)$. \code{theil_sen} is a Python loop of
$m$ fits of up to $200^2/2$ pairs each. \emph{learned} mode adds pSAX's density step, about
$1000\times m$ values per array. No \code{estimate} hook.

\entryfield{Failure modes}
\begin{itemize}
\item \emph{Noise reads as shape.} On pure Gaussian noise (3000 segments) learned mode
  labelled 53.5\% of segments up or down, with 96.6\% of the maximum occupancy entropy; a
  phase-randomised noise floor made 94.6\% of them ``same''
  (\code{Working/Detection/sax/dsax_python/IMPLEMENTATION_NOTES.md}, section 6.3). The
  post-run warning for this (\code{trend_diagnostic_rows}) and the noise-floor estimator
  (\code{surrogate_same_halfwidth}) exist in the code, but the web UI calls neither, and
  the ``Estimate noise floor'' button the \emph{min\_same\_halfwidth} help text mentions is
  not in the web UI.
\item \emph{Drift.} A net trend pushes every segment towards one end of the alphabet
  (section 7.3 of the notes; test 9).
\item \emph{Offset sensitivity.} Shifting the segment grid by a few samples changes many
  letters on signals that turn within a segment: 11.7\% agreement on a triangle wave in the
  notes' section 6.1.
\item \emph{Even alphabets} have a cutline at zero and no ``same'' letter.
\item \emph{Degenerate cutlines} (\code{cutlines_degenerate}) on constant spans, clean
  ramps or colliding cutlines; the explanation exists in \code{degeneracy_explanation} but
  is not shown by the web UI.
\end{itemize}

\entryfield{Validation}
\code{tests/test_dsax_engineered.py} checks the answers on engineered signals: ramps encode
all-up or all-down, a constant is all-same in every mode, a step marks exactly its segment,
a staircase fills a five-letter alphabet in order, a ``sharkfin'' matches its pattern, white
noise stays balanced about zero, drift swamps the encoding and detrending fixes it, the
encoding is invariant to amplitude scaling and to sampling rate at matched seconds, and the
fitted estimators flip less than endpoints under noise. No test checks the strings against
hand-labelled fungal data. The prior art is a family of trend encodings; this one is its
slope channel with a learned zero-anchored alphabet, and claims nothing more.

\entryfield{References}
\citet{esmael2012multivariate} (trend symbols), \citet{sun2014improvement} (the endpoints
trend), \citet{malinowski2013onedsax} (the slope channel of 1d-SAX),
\citet{agrawal1995querying} (up/down/stable shapes), \citet{bountrogiannis2022distribution}
(the quantiser), \citet{theiler1992testing} (the surrogates of the noise-floor estimator).

\entryfield{Tests}
\code{tests/test_dsax.py} (33 tests: segmentation round trip, the delta conversion rule,
modes, details); \code{tests/test_dsax_engineered.py} (17, above);
\code{tests/test_dsax_diagnostics.py} (14, the post-run rows);
\code{tests/test_drop_detection_nan.py} (missing segments, gap-free input unchanged);
\code{tests/test_sax_determinism.py} (fixed seed).
